\section{Classi di complessitá}

\subsection{TIME}
Un linguaggio $L$ appartiene alla classe $\TIME(f(n))$ se esiste una macchina di Turing $M$ che decide $L$ e opera in tempo $f(n)$.
\subsubsection{Teorema: Linear Time Speedup}
Se $L \in \TIME(f(n))$ allora $\forall \varepsilon > 0 \quad L \in \TIME(f'(n))$ \\
Dove $f'(n) = \varepsilon f(n) + n + 2$ \\

\textbf{Prova:} \\

\textbf{QED} \\

Questo teorema ci permette di ignorare costanti moltiplicative nella misurazione della complessitá di tempo.


\subsection{SPACE}
Un linguaggio $L$ appartiene alla classe $\SPACE(f(n))$ se esiste una macchina di Turing $M$ che decide $L$ e opera in spazio $f(n)$.
\subsubsection{Teorema: Linear Space Speedup}
Se $L \in \SPACE(f(n))$ allora $\forall \varepsilon > 0 \quad L \in \SPACE(f'(n))$ \\
Dove $f'(n) = \varepsilon f(n) + n + 2$ \\

\textbf{Prova:} \\

\textbf{QED} \\

Questo teorema ci permette di ignorare costanti moltiplicative nella misurazione della complessitá di spazio.
